% Uses only booktabs (already loaded by appendix_filtering.tex); no extra packages required.

\section{Additional statistical checks}
\label{app:stats}

The main paper reports exact node-permutation tests of the cosine coefficient with
controls for mean and maximum single-vector strength. This appendix gives the full
inferential setup, the supporting tables, and the robustness checks behind that result.
Both schemes analysed here are the normalized injection variants that survive the scheme
filter of Appendix~\ref{app:naive}: normalized-sum (\texttt{normTrue}) and
projection-controlled (\texttt{per\_axis}); both use the eight behaviours that remain after
\texttt{power\_seeking} is dropped (Appendix~\ref{app:naive}, Filter~1).

\paragraph{Why an exact node-permutation null.}
The dyads are not independent. With eight behaviours there are $\binom{8}{2}=28$ pairs, and
each behaviour appears in seven of them, so any two dyads that share a behaviour are
correlated. Edge-wise permutation---shuffling the cosine entries independently---treats the
dyads as exchangeable and so ignores exactly this dependence. We instead use the
node-permutation (MRQAP) null: we draw a permutation $\pi$ of the eight behaviours, apply it
\emph{simultaneously} to the rows and columns of the cosine matrix, and refit, holding the
outcome, the single-vector controls, and the coherence mask fixed at their observed $(i,j)$
positions. Because there are only $8!=40{,}320$ permutations we enumerate all of them: the
resulting $p$-value is \emph{exact}, not Monte-Carlo. Throughout, we regress the standardized
outcome on the standardized mean and maximum single-vector strength and the standardized
cosine $c_{ij}$, and report the partial cosine coefficient $\beta_c$. The two outcomes are the
joint magnitude $M$ and the directional residual $S$; normalized-sum uses all $28$ dyads,
while projection-controlled retains $22$ after the coherence gate.

\paragraph{Bivariate associations.}
The uncontrolled correlations (Table~\ref{tab:app-bivar}) are larger than the controlled
coefficients: for normalized-sum, $r(c_{ij},M)=0.646$ and $r(c_{ij},S)=0.475$; for
projection-controlled, $r(c_{ij},M)=0.642$ and $r(c_{ij},S)=0.472$. These bivariate
relationships are not confirmatory, because they do not control for single-vector strength;
once the mean and maximum single-vector magnitudes enter the model the cosine coefficient
shrinks substantially, and it is the controlled estimate that the node-permutation test
evaluates.

\begin{table}[h]
  \centering
  \small
  \setlength{\tabcolsep}{5pt}
  \begin{tabular}{llcccc}
    \toprule
    scheme & outcome & $\beta_c$ & $p_{\text{node}}$ & $r(c_{ij},\cdot)$ & $p_{\text{biv}}$ \\
    \midrule
    Normalized-sum ($n=28$) & Magnitude $M$ & $+0.188$ & $0.081$ & $0.646$ & $0.0002$ \\
                            & Directional residual $S$ & $+0.310$ & $0.047$ & $0.475$ & $0.011$ \\
    \midrule
    Projection-controlled ($n=22$) & Magnitude $M$ & $+0.187$ & ${<}0.001$ & $0.642$ & $0.001$ \\
                            & Directional residual $S$ & $+0.411$ & ${<}0.001$ & $0.472$ & $0.026$ \\
    \bottomrule
  \end{tabular}
  \caption{Per-scheme dyadic regressions. $\beta_c$ is the standardized partial coefficient
  of $c_{ij}$ after controlling for mean and maximum single-vector strength, with the exact
  node-permutation $p$-value ($p_{\text{node}}$) over all $8!$ behaviour relabellings. The last
  two columns give the uncontrolled bivariate Pearson correlation and its $p$-value.}
  \label{tab:app-bivar}
\end{table}

\paragraph{Edge-wise versus node-permutation inference.}
Table~\ref{tab:app-null} contrasts the two nulls for the same fitted coefficient. For
normalized-sum the two agree to within $0.013$, with the node-permutation $p$ marginally
\emph{more} conservative ($M$: $0.068\!\to\!0.081$; $S$: $0.042\!\to\!0.047$): accounting for
the shared-node dependence does not inflate significance. For projection-controlled the node
permutation instead yields markedly stronger evidence (magnitude $0.148\!\to\!{<}0.001$).
Because the node-permutation null holds the distribution of cosine values fixed and only
relabels behaviours, a small $p$ means the association depends on \emph{which specific
behaviours} are geometrically aligned, not merely on the spread of cosine values---a property
the edge-wise null cannot detect.

\begin{table}[h]
  \centering
  \small
  \setlength{\tabcolsep}{5pt}
  \begin{tabular}{llccc}
    \toprule
    scheme & outcome & $\beta_c$ & $p_{\text{edge}}$ (MC) & $p_{\text{node}}$ ($n!$) \\
    \midrule
    Normalized-sum & Magnitude $M$ & $+0.188$ & $0.068$ & $0.081$ \\
                   & Directional residual $S$ & $+0.310$ & $0.042$ & $0.047$ \\
    \midrule
    Projection-controlled & Magnitude $M$ & $+0.187$ & $0.148$ & ${<}0.001$ \\
                   & Directional residual $S$ & $+0.411$ & $0.047$ & ${<}0.001$ \\
    \bottomrule
  \end{tabular}
  \caption{Edge-wise versus node-permutation $p$-values for the cosine coefficient. The point
  estimate is identical (same fit); only the null differs. Edge-wise uses $10{,}000$
  Monte-Carlo shuffles of the cosine entries; node-permutation is the exact $8!$ relabelling.}
  \label{tab:app-null}
\end{table}

\paragraph{Semantic similarity and signed versus unsigned cosine.}
Adding a semantic-similarity control between behaviour descriptions strengthens rather than
removes the directional cosine coefficient: normalized-sum increases from $+0.310$ to $+0.361$,
and projection-controlled from $+0.411$ to $+0.443$ (Table~\ref{tab:app-dir}). The directional
association is therefore not explained by label-level semantic similarity. The check that most
sharply separates the schemes is replacing signed $c_{ij}$ with its magnitude $|c_{ij}|$: for
normalized-sum the effect collapses ($+0.053$, $p=0.742$), whereas projection-controlled
retains a significant coefficient ($+0.193$, $p<0.001$). The association is thus genuinely
\emph{signed}---the direction of alignment matters, not merely its degree---and this signed
structure is far more pronounced under projection-controlled injection.

\begin{table}[h]
  \centering
  \small
  \setlength{\tabcolsep}{5pt}
  \begin{tabular}{lcccccc}
    \toprule
     & \multicolumn{2}{c}{signed $c_{ij}$} & \multicolumn{2}{c}{$+$ semantic sim.} & \multicolumn{2}{c}{unsigned $|c_{ij}|$} \\
    \cmidrule(lr){2-3}\cmidrule(lr){4-5}\cmidrule(lr){6-7}
    scheme & $\beta$ & $p$ & $\beta$ & $p$ & $\beta$ & $p$ \\
    \midrule
    Normalized-sum        & $+0.310$ & $0.047$    & $+0.361$ & $0.023$    & $+0.053$ & $0.742$    \\
    Projection-controlled & $+0.411$ & ${<}0.001$ & $+0.443$ & ${<}0.001$ & $+0.193$ & ${<}0.001$ \\
    \bottomrule
  \end{tabular}
  \caption{Robustness of the directional ($S$) cosine coefficient under the node-permutation
  null: signed-cosine baseline; with an added semantic-similarity control; and with unsigned
  $|c_{ij}|$ substituted for signed cosine. All $\beta$ are standardized.}
  \label{tab:app-dir}
\end{table}

\paragraph{Leave-one-behaviour-out.}
Each fold drops one behaviour, shrinking the cosine matrix to $7\times7$ and making the
node-permutation null exact at $7!=5{,}040$ (Table~\ref{tab:app-loo}). For normalized-sum every
fold has exactly $\binom{7}{2}=21$ dyads; for projection-controlled the fold size varies
($15$--$19$) because removing a behaviour also changes which remaining dyads pass the coherence
gate. The normalized-sum effect is sensitive to the behaviour set: dropping \emph{apathetic}
reduces the coefficient to $+0.064$ ($p=0.754$), while several other drops preserve or amplify a
positive coefficient (dropping \emph{hallucinating} gives $+0.632$). Projection-controlled
estimates remain positive and significant across all folds, ranging from $+0.246$ to $+0.549$,
though this scheme is also affected by coherence censoring. These checks support the main
interpretation as a moderate, behaviour-set-dependent tendency rather than a sharp rule.

\begin{table}[h]
  \centering
  \small
  \setlength{\tabcolsep}{5pt}
  \begin{tabular}{lccc@{\hskip 2em}ccc}
    \toprule
     & \multicolumn{3}{c}{Normalized-sum} & \multicolumn{3}{c}{Projection-controlled} \\
    \cmidrule(lr){2-4}\cmidrule(lr){5-7}
    dropped behaviour & $n$ & $\beta_c$ & $p_{\text{node}}$ & $n$ & $\beta_c$ & $p_{\text{node}}$ \\
    \midrule
    (none)        & 28 & $+0.310$ & $0.047$    & 22 & $+0.411$ & ${<}0.001$ \\
    \midrule
    apathetic     & 21 & $+0.064$ & $0.754$    & 16 & $+0.246$ & ${<}0.001$ \\
    confidence    & 21 & $+0.405$ & $0.029$    & 16 & $+0.519$ & ${<}0.001$ \\
    evil          & 21 & $+0.418$ & $0.013$    & 16 & $+0.463$ & ${<}0.001$ \\
    formality     & 21 & $+0.263$ & $0.242$    & 16 & $+0.384$ & ${<}0.001$ \\
    hallucinating & 21 & $+0.632$ & ${<}0.001$ & 18 & $+0.549$ & ${<}0.001$ \\
    humorous      & 21 & $+0.288$ & $0.135$    & 19 & $+0.382$ & $0.002$    \\
    impolite      & 21 & $+0.335$ & $0.088$    & 16 & $+0.269$ & ${<}0.001$ \\
    sycophantic   & 21 & $+0.213$ & $0.216$    & 15 & $+0.430$ & ${<}0.001$ \\
    \bottomrule
  \end{tabular}
  \caption{Leave-one-behaviour-out for the directional ($S$) coefficient. $n$ is the number of
  dyads surviving the coherence gate in that fold; $p_{\text{node}}$ is the exact $7!$
  node-permutation $p$-value.}
  \label{tab:app-loo}
\end{table}

\paragraph{Summary.}
\textbf{Projection-controlled (\texttt{per\_axis}).} A consistent, signed, semantically robust,
and fold-stable directional association between behaviour-vector cosine and the steering
outcome, with magnitude also significant under the exact null. \textbf{Normalized-sum
(\texttt{normTrue}).} The same qualitative pattern---positive, signed, and reinforced by the
semantic control---but weaker and dependent on the behaviour set, with the magnitude
coefficient and several leave-one-behaviour-out folds not reaching significance. We therefore
read the geometry--control relationship as a moderate, scheme- and behaviour-set-dependent
tendency rather than a universal law.
